\documentclass{amsart}
% For dealing with references we use the comment environment
\usepackage{verbatim}
\newenvironment{reference}{\comment}{\endcomment}
%\newenvironment{reference}{}{}
\newenvironment{slogan}{\comment}{\endcomment}
\newenvironment{history}{\comment}{\endcomment}

% For commutative diagrams we use Xy-pic
\usepackage[all]{xy}

% We use 2cell for 2-commutative diagrams.
\xyoption{2cell}
\UseAllTwocells

% We use multicol for the list of chapters between chapters
\usepackage{multicol}

% This is generally recommended for better output
\usepackage{lmodern}
\usepackage[T1]{fontenc}

% For cross-file-references

% Package for hypertext links:
\usepackage{hyperref}

% For any local file, say "hello.tex" you want to link to please

% Theorem environments.
%
\theoremstyle{plain}
\newtheorem{theorem}[subsection]{Theorem}
\newtheorem{proposition}[subsection]{Proposition}
\newtheorem{lemma}[subsection]{Lemma}

\theoremstyle{definition}
\newtheorem{definition}[subsection]{Definition}
\newtheorem{example}[subsection]{Example}
\newtheorem{exercise}[subsection]{Exercise}
\newtheorem{situation}[subsection]{Situation}

\theoremstyle{remark}
\newtheorem{remark}[subsection]{Remark}
\newtheorem{remarks}[subsection]{Remarks}

\numberwithin{equation}{subsection}

% Macros
%
\def\lim{\mathop{\mathrm{lim}}\nolimits}
\def\colim{\mathop{\mathrm{colim}}\nolimits}
\def\Spec{\mathop{\mathrm{Spec}}}
\def\Hom{\mathop{\mathrm{Hom}}\nolimits}
\def\Ext{\mathop{\mathrm{Ext}}\nolimits}
\def\SheafHom{\mathop{\mathcal{H}\!\mathit{om}}\nolimits}
\def\SheafExt{\mathop{\mathcal{E}\!\mathit{xt}}\nolimits}
\def\Sch{\mathit{Sch}}
\def\Mor{\mathop{\mathrm{Mor}}\nolimits}
\def\Ob{\mathop{\mathrm{Ob}}\nolimits}
\def\Sh{\mathop{\mathit{Sh}}\nolimits}
\def\NL{\mathop{N\!L}\nolimits}
\def\CH{\mathop{\mathrm{CH}}\nolimits}
\def\proetale{{pro\text{-}\acute{e}tale}}
\def\etale{{\acute{e}tale}}
\def\QCoh{\mathit{QCoh}}
\def\Ker{\mathop{\mathrm{Ker}}}
\def\Im{\mathop{\mathrm{Im}}}
\def\Coker{\mathop{\mathrm{Coker}}}
\def\Coim{\mathop{\mathrm{Coim}}}

% Boxtimes
%
\DeclareMathSymbol{\boxtimes}{\mathbin}{AMSa}{"02}

%
% Macros for moduli stacks/spaces
%
\def\QCohstack{\mathcal{QC}\!\mathit{oh}}
\def\Cohstack{\mathcal{C}\!\mathit{oh}}
\def\Spacesstack{\mathcal{S}\!\mathit{paces}}
\def\Quotfunctor{\mathrm{Quot}}
\def\Hilbfunctor{\mathrm{Hilb}}
\def\Curvesstack{\mathcal{C}\!\mathit{urves}}
\def\Polarizedstack{\mathcal{P}\!\mathit{olarized}}
\def\Complexesstack{\mathcal{C}\!\mathit{omplexes}}
% \Pic is the operator that assigns to X its picard group, usage \Pic(X)
% \Picardstack_{X/B} denotes the Picard stack of X over B
% \Picardfunctor_{X/B} denotes the Picard functor of X over B
\def\Pic{\mathop{\mathrm{Pic}}\nolimits}
\def\Picardstack{\mathcal{P}\!\mathit{ic}}
\def\Picardfunctor{\mathrm{Pic}}
\def\Deformationcategory{\mathcal{D}\!\mathit{ef}}

\setlength{\emergencystretch}{3em}
\begin{document}
\title[Stacks Project, Tag 03GW]{478 --- The quasi-finite locus of a separated finite-type morphism is an open subscheme of its relative normalization}
\maketitle
\noindent Copyright (C) 2005--2025 Johan de Jong. Stacks Project authors. Source: \href{https://stacks.math.columbia.edu/tag/03GW}{Tag 03GW}.

\noindent Editorial difficulty note (not author text). Difficulty H2: The proof must identify the whole inverse image of an open normalization neighbourhood, not merely embed the quasi-finite locus. Elementary étale splitting isolates a finite component, normalization respects the resulting disjoint union, and descent of the local isomorphism completes the global identification..

\noindent Editorial provenance note (not author text). History: excerpt from revision \texttt{a0444\allowbreak 6e57e\allowbreak c1fbc\allowbreak 252a8\allowbreak 71afc\allowbreak ec775\allowbreak 2fb28\allowbreak 07b14}, more-morphisms.tex, lines 12327--12423. Reference presentation was adapted; original-author context and core support blocks were retained or added. The mathematical wording is unchanged. Licensed under GNU FDL 1.2 or later; no invariant sections or cover texts. See ../licenses/COPYING. Context blocks reproduce author definitions and supporting results; they are not additional counted problems.

\begin{definition}
\label{morphisms-definition-normalization-X-in-Y}
Let $f : Y \to X$ be a quasi-compact and quasi-separated morphism of schemes.
Let $\mathcal{O}'$ be the integral closure of $\mathcal{O}_X$ in
$f_*\mathcal{O}_Y$. The {\it normalization of $X$ in $Y$} is the
scheme\footnote{The scheme $X'$ need not be normal, for example if
$Y = X$ and $f = \text{id}_X$, then $X' = X$.}
$$
\nu : X' = \underline{\Spec}_X(\mathcal{O}') \to X
$$
over $X$. It comes equipped with a natural factorization
$$
Y \xrightarrow{f'} X' \xrightarrow{\nu} X
$$
of the initial morphism $f$.
\end{definition}

\begin{lemma}
\label{lemma-finite-type-separated}
Let $f : X \to S$ be a morphism of schemes.
Assume $f$ is of finite type and separated.
Let $S'$ be the normalization of $S$ in $X$, see
Morphisms, Definition \ref{morphisms-definition-normalization-X-in-Y}.
Picture:
$$
\xymatrix{
X \ar[rd]_f \ar[rr]_{f'} & & S' \ar[ld]^\nu \\
& S &
}
$$
Then there exists an open subscheme $U' \subset S'$ such that
\begin{enumerate}
\item $(f')^{-1}(U') \to U'$ is an isomorphism, and
\item $(f')^{-1}(U') \subset X$ is the set of points at which
$f$ is quasi-finite.
\end{enumerate}
\end{lemma}

\begin{proof}
By Morphisms, Lemma \href{https://stacks.math.columbia.edu/tag/01TI}{lemma quasi finite points open (Tag 01TI)}
the subset $U \subset X$ of points where $f$ is quasi-finite is open.
The lemma is equivalent to
\begin{enumerate}
\item[(a)] $U' = f'(U) \subset S'$ is open,
\item[(b)] $U = (f')^{-1}(U')$, and
\item[(c)] $U \to U'$ is an isomorphism.
\end{enumerate}
Let $x \in U$ be arbitrary. We claim there exists an open
neighbourhood $f'(x) \in V \subset S'$ such that $(f')^{-1}V \to V$ is an
isomorphism. We first prove the claim implies the lemma.
Namely, then $(f')^{-1}V \cong V$ is both locally of finite
type over $S$ (as an open subscheme of $X$) and for $v \in V$ the residue
field extension $\kappa(v)/\kappa(\nu(v))$ is algebraic (as
$V \subset S'$ and $S'$ is integral over $S$). Hence the fibres
of $V \to S$ are discrete (Morphisms, Lemma
\href{https://stacks.math.columbia.edu/tag/01TE}{lemma algebraic residue field extension closed point fibre (Tag 01TE)})
and $(f')^{-1}V \to S$ is locally quasi-finite
(Morphisms, Lemma \href{https://stacks.math.columbia.edu/tag/06RT}{lemma locally quasi finite fibres (Tag 06RT)}).
This implies $(f')^{-1}V \subset U$ and $V \subset U'$. Since $x$ was
arbitrary we see that (a), (b), and (c) are true.

\medskip\noindent
Let $s = f(x)$. Let $(T, t) \to (S, s)$ be an elementary \'etale
neighbourhood. Denote by a subscript ${}_T$ the base change to $T$.
Let $y = (x, t) \in X_T$ be the unique point in
the fibre $X_t$ lying over $x$. Note that $U_T \subset X_T$
is the set of points where $f_T$ is quasi-finite, see
Morphisms, Lemma \href{https://stacks.math.columbia.edu/tag/01TM}{lemma base change quasi finite (Tag 01TM)}.
Note that
$$
X_T \xrightarrow{f'_T} S'_T \xrightarrow{\nu_T} T
$$
is the normalization of $T$ in $X_T$, see
Lemma \href{https://stacks.math.columbia.edu/tag/03GV}{lemma normalization smooth localization (Tag 03GV)}.
Suppose that the claim holds for $y \in U_T \subset X_T \to S'_T \to T$, i.e.,
suppose that we can find an open neighbourhood
$f'_T(y) \in V' \subset S'_T$ such that $(f'_T)^{-1}V' \to V'$ is an
isomorphism. The morphism $S'_T \to S'$ is \'etale hence the image
$V \subset S'$ of $V'$ is open. Observe that $f'(x) \in V$ as $f'_T(y) \in V'$.
Observe that
$$
\xymatrix{
(f'_T)^{-1}V' \ar[r] \ar[d] & (f')^{-1}(V) \ar[d] \\
V' \ar[r] & V
}
$$
is a fibre square (as $S'_T \times_{S'} X = X_T$).
Since the left vertical arrow is an isomorphism
and $\{V' \to V\}$ is a \'etale covering, we conclude that the right vertical
arrow is an isomorphism by
Descent, Lemma \href{https://stacks.math.columbia.edu/tag/02L4}{lemma descending property isomorphism (Tag 02L4)}.
In other words, the claim holds for $x \in U \subset X \to S' \to S$.

\medskip\noindent
By the result of the previous paragraph we may replace $S$ by an
elementary \'etale neighbourhood of $s = f(x)$ in order to prove the claim.
Thus we may assume there is a decomposition
$$
X = V \amalg W
$$
into open and closed subschemes where $V \to S$ is finite and $x \in V$,
see Lemma \href{https://stacks.math.columbia.edu/tag/02LN}{lemma etale splits off quasi finite part technical (Tag 02LN)}.
Since $X$ is a disjoint union of $V$ and $W$ over $S$ and since
$V \to S$ is finite we see that the
normalization of $S$ in $X$ is the morphism
$$
X = V \amalg W \longrightarrow V \amalg W' \longrightarrow S
$$
where $W'$ is the normalization of $S$ in $W$, see
Morphisms, Lemmas \href{https://stacks.math.columbia.edu/tag/03GO}{lemma normalization in disjoint union (Tag 03GO)},
\href{https://stacks.math.columbia.edu/tag/01WJ}{lemma finite integral (Tag 01WJ)}, and
\href{https://stacks.math.columbia.edu/tag/03GP}{lemma normalization in integral (Tag 03GP)}.
The claim follows and we win.
\end{proof}
\end{document}
